\section{Fixed domains, random domains, and Walsh coordinates}
\label{mart:section}

The finite-memory weights admit a second explanation: count the Walsh
characters whose last active coordinate was updated at each preceding
layer. We give that derivation first, then compare two ways of obtaining
the opposite-half balance used in Section~\ref{sec:first-route}. A fixed
domain acquires balance from the preceding switch layer; a random domain
also has balance by sampling without replacement. The resulting table
records the quantitative inputs used by several distinct Fourier
transfers. Random sampling is needed for the smaller-subcube argument of
Section~\ref{sec:noise-extensions}, but it is not needed to obtain the
half-density bounds in this section.

We retain $V=\{0,1\}^d$, $n=2^d$, and $G=\Sym(V)$. In the cyclic
coordinate frame, layer $s$ uses $i_s=1+(s\bmod d)$. The position
permutation in that frame is $\Pi_s$. A sweep consists of $d$ layers;
at every sweep boundary the frame is the physical position frame.

\subsection{Walsh coordinates of the adapted flow}
\label{mart:flow-section}

Fix an ordered list $(x_1,\ldots,x_k)$ of distinct input positions and expose the paths of
its first $\ell$ cards.  The list may be deterministic or sampled
independently of the shuffle.  The next card is the tagged card.  Let
$\mathcal F_s$ contain the list and the exposed paths through time $s$,
and let $V_s$ be the set of positions not occupied by exposed cards.
Its size is $m=n-\ell$.  Given a terminal time $T$, define
\[
 p_s(x)=\mathbb P(\Pi_s(x_{\ell+1})=x\mid\mathcal F_T),\qquad
 f_s=p_s-m^{-1}\mathbf1_{V_s}\quad(0\le s\le T).
\]
These are the available set and centered vector of
Section~\ref{sec:first-route}. Lemma~\ref{lem:marginal-averaging}
applies with the exposed cards as blockers: although $p_s$ conditions on
complete paths, it has an $\mathcal F_s$-measurable version. Thus $f_s$
has sum zero, is supported on $V_s$, and has squared norm at most one.

Use the coordinate average $P_i$ from Section~\ref{sec:first-route},
and write $x^{(i)}=x\oplus e_i$ for the partner of $x$.  In the following
expectations we sample the shuffle and, when applicable, the initial
list; we do not fix its terminal exposed paths.  Set
\[
 a_s=\mathbb E\|f_s\|_2^2,\qquad
 L_s=\sum_x f_s(x)^2\mathbf1_{\{x^{(i_s)}\notin V_s\}},
 \qquad w_s=\mathbb E L_s.
\]

With this notation, Lemma~\ref{noise:memory} reads
\begin{equation}\label{mart:noise-update}
 f_{s+1}=P_{i_s}f_s+\xi_s,\qquad
 \mathbb E(\xi_s\mid\mathcal F_s)=0.
\end{equation}
On each edge with exactly one available position $x$, its noise is
$\pm f_s(x)(\delta_x-\delta_{x^{(i_s)}})/2$, with independent fair
signs on distinct such edges.  All other edges contribute zero noise.
For every $t\ge d$,
\begin{equation}\label{mart:energy}
 a_t=\sum_{j=1}^d2^{-j}w_{t-j}.
\end{equation}
Equivalently, with $D_s=\mathbb E\|\xi_s\|_2^2=w_s/2$,
$a_t=\sum_{j=1}^d2^{-(j-1)}D_{t-j}$.
There is also a coordinate-frequency proof, using the finite Boolean version of the
Walsh character basis~\cite{walsh1923}. It avoids expanding the
vector recursion and identifies the same weights by counting the last
update of each character.

\begin{lemma}[Walsh derivation]\label{mart:walsh}
For $A\subseteq\{1,\ldots,d\}$, define
$\widehat f_s(A)=\sum_xf_s(x)(-1)^{\sum_{i\in A}x_i}$.
If $i_s\notin A$, then $\widehat f_{s+1}(A)=\widehat f_s(A)$.
If $i_s\in A$, then
\begin{equation}\label{mart:walsh-update}
 \mathbb E\bigl(\widehat f_{s+1}(A)^2\mid\mathcal F_s\bigr)=L_s.
\end{equation}
These identities imply \eqref{mart:energy}.
\end{lemma}
\begin{proof}
A character omitting the active coordinate is constant on every
updated edge, and the update preserves the sum on that edge.  A
character containing the active coordinate has opposite values at
the endpoints.  An edge with two available positions contributes zero after averaging.
Each edge with one available position contributes its entry times an
independent centered sign. The conditional second moment is therefore $L_s$.

At time $t\ge d$, every nonempty character has a last update among the
previous $d$ layers.  Exactly $2^{d-j}$ subsets have their last update
at time $t-j$: they include that coordinate, omit the $j-1$ more
recent coordinates, and choose arbitrarily among the other $d-j$.
The empty coefficient is zero.  Parseval, with normalization
$\|f\|_2^2=n^{-1}\sum_A\widehat f(A)^2$, now gives
\eqref{mart:energy}.
\end{proof}

\subsection{Opposite halves and large marginals}
\label{mart:balance-section}

The remaining input to the scalar recurrence is a gain when an available
position meets another available position.  The masses and the available set
are generally correlated.  The needed independence holds across the
two halves during the $d-1$ intervening coordinate updates.

\begin{lemma}[Half-space factorization]\label{mart:halves}
Fix $s\ge d-1$, put $a=s-d+1$, and let $H_0,H_1$ be the two halves
defined by coordinate $i_s$.  Conditional on $\mathcal F_a$, for
$h\in\{0,1\}$,
\begin{equation}\label{mart:half-factorization}
 \mathbb E\!\left[
  \sum_{x\in H_h}f_s(x)^2\mathbf1_{\{x^{(i_s)}\in V_s\}}
       \,\middle|\,\mathcal F_a\right]
 =\frac{|V_a\cap H_{1-h}|}{n/2}
   \mathbb E\!\left[\sum_{x\in H_h}f_s(x)^2
       \,\middle|\,\mathcal F_a\right].
\end{equation}
If both half densities at time $a$ are at least $\beta$ except on an
$\mathcal F_a$-measurable event of probability at most $\eta$, then
\begin{equation}\label{mart:noise-bound}
 w_s\le(1-\beta)a_s+\eta.
\end{equation}
\end{lemma}
\begin{proof}
The updates from $a$ through $s-1$ visit each coordinate other than
$i_s$ once and never cross the two halves.  Given the initial history,
the future switch arrays in the halves are independent of that history
and of each other.  Starting from the now known restrictions of $V_a$
and $f_a$, both the exposed paths and the forward flow inside one half
are measurable functions of that half's switch array alone.  The
cylinder description in Lemma~\ref{lem:marginal-averaging} leaves
its unvisited switches independent, so computing the conditional flow
introduces no dependence on the other half.  The flow
value at $x$ uses only the array in its own half, whereas the opposite
partner's membership in $V_s$ uses the other array.  In that other
half, each exposed card is marginally uniform after these $d-1$
layers, since each remaining bit has been refreshed with a fair bit.
Distinctness of the card positions gives the stated expected available
density.  This proves \eqref{mart:half-factorization}.  Subtract the
available-partner contribution from $a_s$ and use $\|f_s\|_2^2\le1$
on the exceptional event to obtain \eqref{mart:noise-bound}.
\end{proof}

We record the two sources of balance, as they yield different
uniformity statements about the initial domain.

\begin{lemma}[Fixed and random initial domains]\label{mart:balance}
Suppose $m\ge n/q$, where $q\ge1$ is a fixed integer.
\begin{enumerate}
\item For every deterministic initial list, \eqref{mart:noise-bound}
holds for $s\ge d$ with
$\beta=1/(2q)$ and $\eta=2\exp(-n/(8q))$.
\item If the initial list is uniformly random and independent of the
shuffle, it holds for $s\ge d-1$ with $\beta=1/(2q)$ and either of
the valid bounds
\[
 \eta=2\exp(-n/(32q)),\qquad
 \eta=2(n+1)\exp\!\left(-\frac{(1-\log2)n}{4q}\right).
\]
\item In the random-list setting it also holds with
$\beta=1/(4q)$ and $\eta=2\exp(-c n/q)$ for an absolute $c>0$,
using the weaker requirement $|V_a\cap H_h|\ge m/8$.
\end{enumerate}
\end{lemma}
\begin{proof}
For the first assertion, $a=s-d+1\ge1$ is the time immediately after
the preceding update in coordinate $i_s$.  Given the history before
that update, either half's available count is $A+\operatorname{Bin}(L,1/2)$,
where $2A+L=m$ and $L\le m$.  A deviation below $m/4$ has probability
at most $\exp(-m/8)$ by the independent bounded-difference estimate
for the $L$ fair choices.
The union bound handles both halves.  Each
half with at least $m/4$ available positions has available density at least $1/(2q)$.

For a random list, $V_a$ is unconditionally a uniform $m$-subset,
because a fixed permutation maps a uniform subset to a uniform subset.
This statement is unconditional; no independence of $f_a$ and $V_a$
is asserted.  A half count has mean $m/2$.  Its Doob martingale under
sequential sampling without replacement has increments of absolute
value at most one: changing the next draw can be coupled with the
remaining completion by exchanging those two sites.  The bounded-increment martingale inequality gives probability at most
$\exp(-m/32)$ for a
downward deviation of $m/4$, proving the first displayed choice.

For the other choice, sample independent Bernoulli variables with
parameter $m/n$ and condition their total to equal $m$.  The latter
event has probability at least $1/(n+1)$ since $m$ is a binomial mode;
the case $m=n$ is deterministic.  Before conditioning, a half count
$X$ satisfies $\mathbb E2^{-X}\le e^{-m/4}$.  Thus
$\mathbb P(X<m/4)\le\exp(-(1-\log2)m/4)$.  Condition and sum over
the two halves.

Finally, failure of the $m/8$ balance forces at least
$k=\lceil7m/8\rceil$ of the $m$ sampled sites into one half.  For any
specified $k$ sample indices the probability is at most $2^{-k}$.
The union bound gives $2\binom mk2^{-k}$, which is at most
$2e^{-cm}$: use $\binom m{m-k}\le(8e)^{m/8}$ and
$(7/8)\log2-(1/8)\log(8e)>0$.  Lemma~\ref{mart:halves}
now proves all three assertions.
\end{proof}

\begin{lemma}[Scalar decay and sequential comparison]\label{mart:decay}
If \eqref{mart:noise-bound} holds for every $s\ge d$, then for any
$\gamma\in[1-\beta/2,1)$,
\begin{equation}\label{mart:decay-bound}
 a_t\le\gamma^{t-2d}+\eta/\beta\qquad(t\ge0).
\end{equation}
At time $T$, the total-variation distance of the images of a list of
$k$ cards from the uniform injection is at most the sum of the
expected conditional tagged distances.  Each such summand is at most
$\frac12\sqrt{ma_T}\le\frac12\sqrt{na_T}$.
The estimates hold for each deterministic list under the first balance
assertion, and on average over lists under the other assertions.
\end{lemma}
\begin{proof}
Apply \eqref{noise:comparison-solution} with
$\rho=1-\beta$, $s_0=d$, $t_0=2d$, and $a_0\le1$.
The condition on $\gamma$ is exactly
$(1-\beta)/(2\gamma-1)\le1$.
For each list entry, Cauchy--Schwarz bounds the mean path-conditioned
variation by $\frac12\E\|f_T\|_1\le\frac12\sqrt{ma_T}$.
Lemma~\ref{lem:sequential-tv} sums these errors: its uniform reference
depends only on the previously exposed endpoints, so its path-conditioning
version applies. Averaging this inequality over an independently sampled
list gives the random-domain assertion.
\end{proof}

\begin{proposition}[Quantitative marginal estimates]\label{mart:marginals}
The following choices give the indicated energy bounds uniformly over
the sequential indices, and hence the corresponding large-marginal
estimates.  In each row $q$ is fixed and at most $n-n/q$ cards are
observed.  The letters $C_{\mathrm{fix}},L,C,b_0$ denote fixed integers satisfying
the displayed conditions.
\begin{center}
\begin{tabular}{c c c c c}
Input used & $q$ & $\beta$ & $\gamma$ & Time and energy\\ \hline
fixed & $16$ & $1/32$ & $127/128$ &
$T=C_{\mathrm{fix}}d$, $a_T=O(n^{-6})$\\
fixed & $8$ & $1/16$ & $63/64$ &
$T=512d$, $a_T=O(n^{-7})$\\
random & $32$ & $1/64$ & $127/128$ &
$T=Ld$, $a_T=O(n^{-10})$\\
random & $8$ & $1/16$ & $31/32$ &
$T=200d$, $a_T=O(n^{-8})$\\
random & $256$ & $1/1024$ & $2047/2048$ &
$T=Cd$, $a_T=o(n^{-4})$\\
random & $64$ & $1/128$ & $511/512$ &
$T=b_0d$, $a_T=O(n^{-5})$\\
random & $K=8192$ & $1/(2K)$ & $1-1/(8K)$ &
$T=(100K+2)d$, $a_T=o(n^{-4})$
\end{tabular}
\end{center}
The integer choices are
\[
 (127/128)^{C_{\mathrm{fix}}-2}\le2^{-6},\quad
 (127/128)^{L-2}\le2^{-10},\quad
 (2047/2048)^{C-2}<2^{-4},\quad
 (511/512)^{b_0-2}\le2^{-5}.
\]
The first two rows will be used for prescribed lists, and the other
rows on average over uniform initial lists. In fact the fixed-domain
part of Lemma~\ref{mart:balance} also proves every energy bound in the
table for each prescribed list: it has at least the displayed $\beta$
and an exponentially small exceptional term. The distinction in the
table records the information used by the transfers below.
In particular the
eight-block random-domain row has average marginal distance
$O(n^{-5/2})$, and the $64$-block row has distance $O(n^{-1})$.
\end{proposition}
\begin{proof}
Apply Lemmas~\ref{mart:balance} and \ref{mart:decay}.  The $256$-block
row uses the $m/8$ threshold; the other rows use $m/4$.  The exceptional
terms are exponentially small in $n$.  For the numerical choices,
$(63/64)^{510d}=O(2^{-7d})$ and
$(31/32)^{198d}\le2^{-8d}$, since
$198\log_2(32/31)>8$.  In the last row the geometric term is at most
$\exp(-100d/8)=o(2^{-4d})$.  The sequential estimate multiplies
$\sqrt{na_T}/2$ by at most $n$, giving the claimed errors.
\end{proof}


\subsection{Using the marginal bounds on the full permutation group}
\label{mart:applications}

We now specify how each row of Proposition~\ref{mart:marginals} supplies
a full-deck bound. This passage uses the abstract results of
\emph{From partial permutation information to Fourier bounds}
\cite{partialFourier}. We state the needed bounds at their point of use.
They use the mass transform
$\widehat\nu(\rho)=\sum_{g\in G}\nu(g)\rho(g)$ and unnormalized
Hilbert--Schmidt norm. For $u>0$, put
$C_u=\sup_{m\ge1}\sum_{\tau\in\widehat{S_m}}(\dim\tau)^{-u}$.
Theorem~\ref*{l-l:degree-all-powers} of~\cite{partialFourier} gives $C_u<\infty$
and $\sum_{\rho\ne\mathbf1,\sgn}(\dim\rho)^{-u}=o(1)$ as $n\to\infty$.

For a partition into $q$ equal blocks, retain the subgroups
$H_i=\Sym(D_i)$ and left-coset convention of
Section~\ref{sec:probability-completion}. A fixed-list bound supplies
every such partition; an averaged bound supplies one by averaging
the sum of its $q$ complementary marginal errors. In either case write
that sum as $\delta_n=o(1)$.

We use the common trimming operations from
\cite[Lemma~\ref*{l-l:trim}]{partialFourier}. Heavy-coset deletion and
normalization give a probability within $o(1)$ of the block law,
with cap $B/[G:H_i]$ for $B=3$ or $4$. Successive excess trimming
instead gives a subprobability of mass at least $1-\delta_n$ and cap
one. Heavy-coset deletion without normalization gives mass
$1-o(1)$ and cap two. Each operation produces one measure satisfying
all the caps simultaneously.

\begin{proposition}[Quantitative applications of the martingale bounds]
\label{mart:application-times}
Use the integer parameters specified in Proposition~\ref{mart:marginals}.
Every total time in Table~\ref{mart:transfer-table} gives worst-start
total-variation convergence to uniform as $d\to\infty$.
\end{proposition}

The table distinguishes a bound on the operator norm from one on the
\emph{squared}, unnormalized Hilbert--Schmidt norm. For its first three
rows using squared Hilbert--Schmidt bounds, the exact inequalities are
\begin{align}
 D\|\widehat\nu(\rho)\|_{\HS}^2
 &\le9qC_2D^{4/q} &&(q=16,\ B=3),\label{mart:hs16}\\
 \|\widehat\nu(\rho)\|_{\HS}^2
 &\le qC_2D^{-1+6/q} &&(q=8,\ B=1),\label{mart:hs8}\\
 \|\widehat\nu(\rho)\|_{\HS}^2
 &\le4qC_{16}D^{-1+18/q} &&(q=64,\ B=4).\label{mart:hs64}
\end{align}
Here $\nu$ denotes the retained measure specified in the corresponding
row, and $D=\dim\rho$. The constants in the other rows are also
independent of $n$, except for $M=1-o(1)$, the retained probability
in the random-partition construction.

\begingroup
\small
\setlength{\tabcolsep}{4pt}
\renewcommand{\arraystretch}{1.2}
\begin{longtable}{@{}p{.31\linewidth}p{.30\linewidth}r r r@{}}
\caption{Fourier substitutions for the marginal bounds. The column $r$
counts independent blocks; $\kappa$ bounds the exponent of $D$ in
their nonsign Plancherel summand, up to a fixed factor.}
\label{mart:transfer-table}\\
\toprule
Transfer and retained law & Fourier bound & $r$ & $\kappa$ & Total time$/d$\\
\midrule
\endfirsthead
\toprule
Transfer and retained law & Fourier bound & $r$ & $\kappa$ & Total time$/d$\\
\midrule
\endhead
\bottomrule
\endfoot
Random partition, $q=8$;\newline
probability, mass $M$ before conditioning
&
$\op:\ \dfrac2M\sqrt{16C_1}D^{-5/16}$
& $8$ & $-3$ & $1600$\\
Coefficient space, $q=16$;\newline probability, cap $3$
&
HS squared: \eqref{mart:hs16}
& $4$ & $-2$ & $4C_{\mathrm{fix}}$\\
Pointwise character, $q=8$;\newline subprobability, cap $1$
&
HS squared: \eqref{mart:hs8}
& $16$ & $-3$ & $8192$\\
Global isotypic, $q=64$;\newline probability, cap $4$
&
HS squared: \eqref{mart:hs64}
& $32$ & $-22$ & $32b_0$\\
Coefficient average, $q=32$;\newline probability, cap $4$
&
$\op:\ \dfrac83\sqrt{C_1}D^{-5/16}$
& $8$ & $-3$ & $8L$\\
Orbit average, $q=256$;\newline probability, cap $4$
&
$\op:\ 4\sqrt{C_{32}}D^{-1/4}$
& $72$ & $-34$ & $72C$\\
Small-index character,\newline $K=q=8192$;\newline subprobability, cap $2$
&
$\HS^2:\ 2KC_1D^{-1+1026/K}$
& $4$ & $\le-1$ & $3276808$\\
\end{longtable}
\endgroup

\begin{proof}
All Fourier bounds in the table are exact cases of
\cite{partialFourier}: Proposition~\ref*{l-l:random-partition} for
the first row; Proposition~\ref*{l-l:coefficient-lift} for the
coefficient and pointwise-character rows; Proposition~\ref*{l-l:isotypic}
for the global isotypic row;
Propositions~\ref*{l-l:coefficient-average} and
\ref*{l-l:orbit-average} for the two projection averages; and
Theorem~\ref*{l-l:character-lift} for the last row.
The last transfer uses a one-dimensional character of index at most
$b^{1024}$ in every block type of degree $b$, which gives
$(1024+2)/K=1026/8192<1/4$ in its displayed bound.

The first row uses the averaged complementary-coset estimate at
$200d$ directly. Its transfer constructs a conditional probability
law and keeps the partition random inside the coefficient estimate;
it does not select a deterministic partition there. All other rows
use the selected partition and the indicated common trimming.

For an operator bound $A D^{-\beta}$, the probability completion of
Section~\ref{sec:probability-completion} gives a nonsign summand
$A^{2r}D^{2-2r\beta}$. For a squared Hilbert--Schmidt bound
$A D^{-\beta}$, it gives $A^rD^{1-r\beta}$. Substitution gives the
column $\kappa$. Every displayed upper exponent is negative, so the
corresponding reciprocal-degree sum tends to zero. In particular
the pointwise-character and small-index rows use subprobabilities;
they need the same norm calculation with their mass retained.

To finish those two rows, let $z=\nu(G)=1-o(1)$. Since the original
shuffle block has zero sign expectation, the retained subprobability
has sign coefficient at most $1-z$. Equation~\eqref{noise:subprob-plancherel}
makes $\|\nu^{*r}-z^rU_G\|_1=o(1)$, while restoring the missing mass
costs $1-z^r\le r(1-z)$ in each of the two positive measures.
For the probability rows, closeness to the original block law gives
sign coefficient $o(1)$ and permits replacement of the $r$ factors,
as in Section~\ref{sec:probability-completion}.

Finally multiply the corresponding block length from
Proposition~\ref{mart:marginals} by $r$. For the last row this is
$4(100\cdot8192+2)d=3276808d$. All block lengths are multiples of
$d$, so their convolutions are physical shuffle blocks. Translation
by a deterministic starting deck preserves distance from uniform.
\end{proof}
